\documentclass[11pt]{beamer}
\usepackage[utf8]{inputenc}
\usepackage{amsmath, amssymb, graphicx}
\usepackage{bm} % For bold math symbols (vectors)

% Presentation Theme
\usetheme{Madrid}
\usecolortheme{default}

\title[GMAT \& The Shooting Method]{Astrodynamics 2026-1}
\subtitle{High-Fidelity Targeting: The Differential Corrector \\ \& The Shooting Method}
\author{Prof. Juan}
\institute{Universidad de Antioquia - Aerospace Engineering}
\date{Spring 2026}

\begin{document}

% --- Title Slide ---
\begin{frame}
    \titlepage
\end{frame}

% ==========================================
% SECTION 1: THE HIGH-FIDELITY PROBLEM
% ==========================================
\section{The High-Fidelity Problem}

% --- Slide 1 ---
\begin{frame}{Bridging Theory and Reality}
    We have successfully used Python and Lambert's problem to find the optimal transfer geometry. Why must we now use NASA GMAT?
    
    \vspace{0.3cm}
    \textbf{The Analytical Assumption:} Lambert solvers assume ideal 2-Body physics (only the Sun's point-mass gravity acts on the spacecraft).
    
    \vspace{0.3cm}
    \textbf{The High-Fidelity Reality:} When we ignite our engines in Low Earth Orbit, our spacecraft immediately experiences:
    \begin{itemize}
        \item Earth's Oblateness ($J_2$ perturbations).
        \item Atmospheric Drag.
        \item Solar Radiation Pressure (SRP).
        \item N-Body Gravity (The Moon, Jupiter, etc.).
    \end{itemize}
    \vspace{0.2cm}
    \textit{Because of these perturbations, our perfect analytical guess will always "miss" the target in a real mission.}
\end{frame}

% ==========================================
% SECTION 2: THE SHOOTING METHOD
% ==========================================
\section{The Shooting Method}

% --- Slide 2 ---
\begin{frame}{The Core Concept: The Shooting Method}
    To hit our target in a chaotic physical environment, numerical solvers use a technique called the \textbf{Shooting Method}. 
    
    \vspace{0.3cm}
    \textbf{The Artillery Analogy:}
    \begin{enumerate}
        \item You guess an angle and gunpowder charge, fire the cannon, and watch where it lands.
        \item It lands 10 meters short.
        \item You add a tiny perturbation of gunpowder, fire again, and see that it now lands 8 meters short.
        \item By measuring how the distance changed based on that tiny addition, you calculate the "slope" (the derivative).
        \item You use that derivative to calculate exactly how much gunpowder is needed to eliminate the remaining error.
    \end{enumerate}
\end{frame}

% ==========================================
% SECTION 3: GMAT ARCHITECTURE
% ==========================================
\section{GMAT Architecture}

% --- Slide 3 ---
\begin{frame}{GMAT Architecture: The Target Sequence}
    NASA GMAT implements this exact methodology using three specific script commands inside a \texttt{Target} sequence.
    
    \vspace{0.3cm}
    \begin{itemize}
        \item \textbf{Vary (The Independent Variables):} These are the "knobs" GMAT is allowed to turn (e.g., Burn Magnitude $\Delta V$, True Anomaly). We initialize these with our analytical Lambert guesses.
        \item \textbf{Propagate (The Function):} GMAT runs its high-fidelity numerical integrator (e.g., Runge-Kutta 8/9), moving the spacecraft forward in time subject to all real-world physics.
        \item \textbf{Achieve (The Dependent Variables / Goals):} These are the target conditions we must hit (e.g., Lunar Flyby Radius, Arrival $C_3$, B-Plane coordinates).
    \end{itemize}
\end{frame}

% ==========================================
% SECTION 4: THE MATHEMATICS
% ==========================================
\section{The Mathematics}

% --- Slide 4 ---
\begin{frame}{The Mathematics: The Jacobian Matrix}
    How does GMAT calculate the corrections? It builds a \textbf{Jacobian Matrix ($\mathbf{J}$)}.
    
    \vspace{0.2cm}
    Let $\mathbf{X}$ be the vector of our variables (e.g., $\Delta V$) and $\mathbf{F}(\mathbf{X})$ be the error in our targets. We want $\mathbf{F}(\mathbf{X}) = 0$. GMAT uses forward finite-differences:
    
    \begin{enumerate}
        \item \textbf{Nominal Run:} Propagate the initial guess $\mathbf{X}_0$ and record the error.
        \item \textbf{Perturbation Run:} Add a tiny perturbation (defined in your script) to the first variable. Propagate the entire trajectory again and record the new error.
        \item \textbf{Calculate Gradients:} Calculate the partial derivative. \textit{("If I change $\Delta V$ by 0.0001 km/s, my flyby radius changes by 50 km.")} This populates the Jacobian.
    \end{enumerate}
\end{frame}

% --- Slide 5 ---
\begin{frame}{The Iterative Update}
    Once the Jacobian matrix ($\mathbf{J}$) is fully populated for all variables, GMAT applies a multi-variable Newton-Raphson update.
    
    \vspace{0.3cm}
    It inverts the matrix to find the macroscopic correction required for the next iteration:
    \begin{equation}
        \mathbf{X}_{i+1} = \mathbf{X}_i - \mathbf{J}^{-1} \mathbf{F}(\mathbf{X}_i)
    \end{equation}
    
    \vspace{0.3cm}
    \begin{itemize}
        \item $\mathbf{X}_{i+1}$: The new, updated guesses for the \texttt{Vary} commands.
        \item GMAT repeats this entire process, firing the "cannon" over and over, until the error $\mathbf{F}(\mathbf{X})$ drops below your defined tolerance.
    \end{itemize}
\end{frame}

\end{document}